\documentclass[11pt]{article}
\usepackage[a4paper,margin=1in]{geometry}
\usepackage{fontspec}
\usepackage[boldfont=false]{xeCJK}
\setCJKmainfont{FandolSong-Regular.otf}
\usepackage{amsmath}
\usepackage{amssymb}
\usepackage{array}
\usepackage{booktabs}
\usepackage{graphicx}
\usepackage{adjustbox}
\usepackage{enumitem}
\setlist[itemize]{topsep=2pt,partopsep=0pt,parsep=0pt,itemsep=2pt,leftmargin=1.4em}
\setlist[enumerate]{topsep=2pt,partopsep=0pt,parsep=0pt,itemsep=2pt,leftmargin=1.6em}
\usepackage{parskip}
\usepackage{fvextra}
\fvset{breaklines=true,breakanywhere=true,fontsize=\small}
\setlength{\parindent}{0pt}

\begin{document}

\begin{center}
  {\LARGE\bfseries International trade and globalisation}\\[0.35em]
  {\large Igcse Economics}
\end{center}
\vspace{0.6em}

\section*{Specialisation and free trade}

Countries \textbf{trade} 贸易 with each other: they sell some goods and buy others. Just as one worker can focus on one task, a whole country can focus on making the goods it is best at. This is \textbf{specialisation} 专业化 at a national level.

When a country specialises, it sells the goods it makes well (\textbf{exports} 出口) and buys the goods other countries make well (\textbf{imports} 进口).

\subsection*{Free trade}

\textbf{Free trade} 自由贸易 means trade between countries with no barriers — no taxes or limits on imports. Its benefits:

\begin{itemize}
\item countries can buy some goods more cheaply from abroad than by making them at home.

\item consumers get a wider choice of goods.

\item firms can sell to bigger markets, so they gain economies of scale.

\item more competition makes \textbf{domestic} 国内的 firms (firms in your own country) work harder and improve.

\end{itemize}

All of this can raise \textbf{living standards} 生活水平: lower prices and more goods for everyone.

Free trade also has disadvantages:

\begin{itemize}
\item domestic firms may not be able to \textbf{compete} 竞争 with cheaper foreign goods, so they close. This causes job losses.

\item a country can become too \textbf{dependent} 依赖 on other countries for important goods.

\item new, small industries may never grow if they face strong foreign rivals from the start.

\end{itemize}

\section*{Globalisation and trade protection}

\textbf{Globalisation} 全球化 is the way the world's economies are becoming more and more joined together — through trade, travel, the internet, and firms working across borders.

\subsection*{Multinational companies}

A \textbf{multinational company} 跨国公司 (MNC) is a firm that produces in more than one country — for example, a car maker with factories in several countries. The country an MNC sets up in is the \textbf{host country} 东道国.

\begin{itemize}
\item \textbf{Benefits} to the host country: new jobs, investment, new skills and technology, and more tax.

\item \textbf{Drawbacks}: profits may flow back to the MNC's home country; local firms may be pushed out; and some MNCs pay low wages or pollute (for example, by sending waste to poorer countries).

\end{itemize}

\subsection*{Trade protection}

\textbf{Trade protection} 贸易保护 is a government putting up barriers to limit imports and protect its own firms. The main methods:

\begin{itemize}
\item a \textbf{tariff} 关税 — a tax on imports, which makes them dearer.

\item a \textbf{quota} 配额 — a limit on the \emph{amount} of a good that can be imported.

\item a \textbf{subsidy} 补贴 — money paid to domestic firms so they can sell more cheaply than foreign rivals.

\item an \textbf{embargo} 禁运 — a total ban on trade in a good, or with a country.

\end{itemize}

Reasons for protection: to protect jobs; to protect an \textbf{infant industry} 幼稚产业 (a new industry too small to compete yet); to stop \textbf{dumping} 倾销 (foreign firms selling below cost to destroy local rivals); and to raise government revenue from tariffs.

But protection has costs. Prices rise and choice falls for consumers; protected firms stay weak and lazy; and other countries may hit back with their own barriers — a trade war.

\section*{Foreign exchange rates}

The \textbf{exchange rate} 汇率 is the price of one \textbf{currency} 货币 in terms of another — for example, £1 = \$1.30.

\subsection*{How the rate is set}

In a \textbf{floating exchange rate} 浮动汇率 system, the rate is set by the \textbf{demand} 需求 for and \textbf{supply} 供给 of the currency, just like any market.

\begin{itemize}
\item if more people want a currency — to buy that country's exports, or to invest there — demand for it rises, so its price (the exchange rate) rises.

\item for example, if a country's interest rates rise, foreigners want to save money there, so demand for its currency rises and the rate goes up.

\end{itemize}

\subsection*{Appreciation and depreciation}

\begin{itemize}
\item \textbf{appreciation} 升值 — the currency rises in value (it buys more foreign money).

\item \textbf{depreciation} 贬值 — the currency falls in value.

\end{itemize}

A change in the rate changes the price of exports and imports:

\begin{itemize}
\item if a currency appreciates (rises), its exports become dearer for foreigners and its imports become cheaper. So exports fall and imports rise.

\item if a currency depreciates (falls), its exports become cheaper for foreigners and its imports become dearer. So exports rise and imports fall.

\end{itemize}

\section*{The current account of the balance of payments}

The \textbf{balance of payments} 国际收支 is a record of all the money flowing into and out of a country. Its main part is the \textbf{current account} 经常账户, which has four sections:

\begin{itemize}
\item \textbf{trade in goods} 货物贸易 — exports and imports of physical goods (cars, food, oil).

\item \textbf{trade in services} 服务贸易 — exports and imports of services (tourism, banking, shipping).

\item \textbf{primary income} 初次收入 — wages, interest, and profit earned from abroad.

\item \textbf{secondary income} 二次收入 — money sent with nothing given in return, such as aid and gifts.

\end{itemize}

\subsection*{Deficit and surplus}

\begin{itemize}
\item a \textbf{current account deficit} 逆差 — more money flows out than in (the country buys more from abroad than it sells).

\item a \textbf{current account surplus} 顺差 — more money flows in than out.

\end{itemize}

\subsection*{Correcting a deficit}

A country with a large deficit must borrow money or sell assets to pay for it. To correct a deficit, a government can:

\begin{itemize}
\item cut total demand, so people buy fewer imports.

\item let the currency depreciate, to make exports cheaper and imports dearer.

\item use trade protection (tariffs and quotas) to cut imports.

\item use supply-side policy to improve the \textbf{competitiveness} 竞争力 of its firms, so exports rise.

\end{itemize}


\end{document}
